\section{Results}\label{s:results}

\subsection{Search results}\label{s:search-results}
The initial search across PubMed, ScienceDirect, and Google Scholar yielded a total of 13,002 articles. After removing 376 duplicates, the screening process based on title and abstract, with the help of PubMed [\term{tiab}] terms, excluded 12,542 articles for not meeting the eligibility criteria, mainly due to irrelevance to the research question or study design. A subsequent full-text review of 61 articles led to further exclusion of 51 articles, primarily for not fully satisfying the inclusion criteria. This process culminated in the inclusion of 10 studies in the final systematic review. The PRISMA flowchart (Figure \ref{fig:prisma-flowchart}) provides a detailed overview of the study selection process.

\begin{figure}[H]
    % \includegraphics[width=0.47\textwidth]{PRISMA Flowchart} % Export from .docx to .pdf -> Trimmed Whitespace
    \includegraphics[width=0.47\textwidth]{PRISMA} % Custom draw.io to .pdf export
    \figcap{PRISMA flow diagram showing the study selection process.}\label{fig:prisma-flowchart}
\end{figure}

\subsection{Characteristics of included studies: summary}\label{s:characteristics-summary}
The systematic review included 10 \newterm{\RCTs} with a mean age range of participants from 70.35 to 79.3 years. The sample size of the studies varied from 45 to 1,098 participants, with male participation ranging from 38\% to 53.76\%. The studies investigated various interventions, including \newterm{\TreatAndExtend} regimens with \newterm{ranibizumab}, \newterm{bevacizumab}, \newterm{ketorolac} eyedrops, \newterm{indomethacin}, and different dosages of ranibizumab. Comparators varied from monthly regimens to placebo and alternative anti-VEGF therapies. The studies primarily assessed the outcomes in terms of \newterm{\BCVA}, \newterm{\CMT}, and safety profiles, including the incidence of adverse events.

\vfill

\subsection{Detailed Results of Included Studies}\label{s:detailed-results}
In their randomized controlled trial, \citet{kertes2020} compared the effectiveness of ranibizumab administered via a \TreatAndExtend regimen against monthly dosing in patients with \nAMD over 24 months. The study reported comparable visual improvements between the T\&E and monthly regimens, with the T\&E approach requiring fewer injections and visits. Similarly, \citet{kertes2019} found the T\&E regimen equally effective as monthly dosing in achieving visual outcomes, but with reduced treatment frequency. \citet{wykoff2018} supported these findings, noting favorable outcomes with the TREX approach.

\citet{nunes2019} investigated the efficacy and safety of ranibizumab and bevacizumab in treating exudative AMD, concluding no significant differences between the two in terms of effectiveness or safety. \citet{russo2018} examined the combination of indomethacin eye drops with \IVR against ranibizumab monotherapy in treating \CNV. Their findings suggested that the combination therapy led to superior anatomical and visual outcomes compared to IVR alone, potentially reducing the frequency of IVR injections.

\citet{eldem2015} compared visual outcomes and treatment frequency in CNV secondary to AMD patients, following either a \WaitAndExtend or \TreatAndObserve regimen. The study concluded that the W\&E regimen was as effective and safe as the T\&O regimen, potentially reducing the burden of follow-up visits.

\citet{schmidt2011} evaluated a quarterly versus monthly ranibizumab regimen for subfoveal CNV secondary to AMD. They observed that while the quarterly regimen maintained BCVA, the improvement in BCVA was greater in the monthly regimen at 12 months, challenging the non-inferiority of the quarterly regimen.

\citet{semeraro2015} explored the effectiveness of ketorolac eye drops in combination with IVR or verteporfin photodynamic therapy plus IVR, compared to IVR monotherapy. The study found superior outcomes in BCVA and \newterm{\CRT} with ketorolac and ranibizumab combination, also requiring fewer IVR injections.

\citet{ho2014} assessed the efficacy and safety of \qty{0.5}{\milli\gram} and \qty{2.0}{\milli\gram} ranibizumab, administered monthly or PRN, in nAMD patients over 24 months. All treatment groups showed meaningful improvements in BCVA by month 24, with no new safety concerns identified.

Finally, the \citet{catt2012} reported comparable effects on visual acuity between ranibizumab and bevacizumab over two years, with PRN treatment leading to less improvement compared to monthly treatment, and no significant differences in mortality or arteriothrombotic events.

\subsection{Narrative Summary of Study Findings}\label{s:narrative-summary}
The systematic review encapsulated a spectrum of methodologies and outcomes, reflecting a comprehensive exploration of anti-angiogenic therapies in AMD. The \TreatAndExtend regimens, prominently featured in studies by \citet{kertes2019, kertes2020}, demonstrated a pivotal shift in treatment strategy, offering comparable visual outcomes to monthly dosing while significantly reducing treatment frequency and patient burden. Innovative approaches such as the combination therapies explored by \citet{russo2018} and \citet{semeraro2015} unveiled promising advancements, showing superior structural and visual outcomes, suggesting potential paradigm shifts in therapeutic strategies.

\subsection{Results of data extraction}\label{s:data-extraction-results}
The detailed results of the data extraction are presented in Table~\ref{tab:descriptor-table}.